\documentclass[aspectratio=169]{beamer}

\usetheme{CambridgeUS}
\usecolortheme{dolphin}

\usepackage[T1]{fontenc}
\usepackage{lmodern}
\usepackage{booktabs}
\usepackage{amsmath}
\usepackage{hyperref}

% Title page details
\title{Improving Autoencoder Recommender Performance}
\subtitle{Masked Loss, Denoising, and Stable Validation}
\author{Group 17}
\date{\today}

\begin{document}

\begin{frame}
	\titlepage
\end{frame}

\begin{frame}{Outline}
	\tableofcontents
\end{frame}

\section{Context}

\begin{frame}{Problem Setup}
	\begin{itemize}
		\item We build a client--product interaction matrix $R \in \mathbb{R}^{m\times n}$.
		\item Entries are dense scores derived from need propensity (SVM probabilities) and risk alignment.
		\item We evaluate recommenders with an \textbf{entry-level holdout}: set a subset of non-zero entries to zero in $R_{\text{train}}$.
	\end{itemize}
	\vspace{0.5em}
	Key risk: after masking, zeros in $R_{\text{train}}$ may mean \emph{``unknown''}, not \emph{``true zero''}.
\end{frame}

\begin{frame}{Baseline Autoencoder (What Was Wrong)}
	Baseline training minimized full MSE:
	\[
		\mathcal{L}_{\text{MSE}} = \frac{1}{mn}\sum_{u,i} (\hat{R}_{ui} - R_{ui})^2
	\]
	\vspace{0.5em}
	When using $R_{\text{train}}$ (with holdout entries set to 0):
	\begin{itemize}
		\item The model is \textbf{penalized} for predicting the held-out positives.
		\item It learns to reproduce the artificial zeros, hurting generalization to missing entries.
	\end{itemize}
\end{frame}

\section{Method}

\begin{frame}{Fix 1: Masked Reconstruction Loss}
	Train only on observed (non-zero) entries:
	\[
		\mathcal{L}_{\text{masked}} = \frac{\sum_{u,i} M_{ui}(\hat{R}_{ui}-R_{ui})^2}{\sum_{u,i} M_{ui} + \varepsilon}
	\]
	where
	\[
		M_{ui} = \mathbb{1}[R_{ui} > 0].
	\]
	\vspace{0.5em}
	Effect:
	\begin{itemize}
		\item holdout-masked zeros do not contribute to the loss
		\item model is free to assign positive scores to held-out entries
	\end{itemize}
\end{frame}

\begin{frame}{Fix 2: Denoising Autoencoder}
	During training we randomly drop a fraction of observed inputs:
	\[
		\tilde{R} = R \odot (1 - D \odot M), \quad D_{ui} \sim \text{Bernoulli}(p)
	\]
	\vspace{0.5em}
	Intuition:
	\begin{itemize}
		\item forces the network to infer missing values from patterns across products
		\item reduces identity mapping / memorization
	\end{itemize}
	\vspace{0.5em}
	Used corruption rate: $p=0.15$.
\end{frame}

\begin{frame}{Fix 3: Stable Optimization and Validation}
	\begin{itemize}
		\item Optimizer: AdamW (better decoupled weight decay).
		\item LR scheduling: ReduceLROnPlateau on entry-level validation MSE.
		\item Early stopping: keep best weights; stop after no improvement.
	\end{itemize}
	\vspace{0.5em}
	Validation signal:
	\begin{itemize}
		\item take a small subset of observed entries in $R_{\text{train}}$ as validation pairs
		\item monitor MSE on those entries (not the full matrix)
	\end{itemize}
\end{frame}

\section{Results}

\begin{frame}{Holdout Validation Metrics (Before vs After)}
	\begin{table}
		\centering
		\begin{tabular}{lrrrr}
			\toprule
			Method                              & RMSE       & Precision@3 & Recall@3 & NDCG@3 \\
			\midrule
			Baseline AE (full MSE)              & 0.44       & 0.93        & 0.93     & 0.93   \\
			Improved AE (masked + denoise + ES) & 0.02--0.03 & 0.99        & 0.99     & 1.00   \\
			\bottomrule
		\end{tabular}
	\end{table}
	\vspace{0.5em}
	Notes:
	\begin{itemize}
		\item Metrics computed on the same entry-level holdout procedure.
		\item RMSE improvement is driven mainly by correcting the training objective (masked loss).
	\end{itemize}
\end{frame}

\begin{frame}{Why the Improvement Is So Large}
	The baseline objective was mismatched to the evaluation protocol:
	\begin{itemize}
		\item Evaluation asks: \emph{can we recover held-out positives?}
		\item Baseline training effectively asked: \emph{can we reproduce the masked zeros?}
	\end{itemize}
	\vspace{0.5em}
	Masked loss aligns training with the intended meaning of zeros:
	\begin{itemize}
		\item zero can mean \textbf{missing/unobserved} (do not train on it)
		\item rather than a true negative preference
	\end{itemize}
\end{frame}

\section{Implementation}

\begin{frame}{What Changed in Code (Notebook)}
	In \texttt{recommend-system.ipynb}:
	\begin{itemize}
		\item Added \texttt{masked\_mse\_loss(pred, target, mask)}.
		\item Training: AdamW + denoising input corruption.
		\item Validation loop: masked loss + entry-level early stopping + LR schedule.
	\end{itemize}
	\vspace{0.5em}
	Tip: keep the architecture small (only 11 products) and focus on correct objective + regularization.
\end{frame}

\begin{frame}{Next Steps}
	\begin{itemize}
		\item Add a stronger recommender metric: hit-rate / NDCG over all items (not only holdout item subsets).
		\item Compare AE vs SVD under the same evaluation definition.
		\item Try a VAE-style objective (optional) if we need calibrated uncertainty.
	\end{itemize}
\end{frame}

\end{document}
